\section{\ComputationCircuits{}}

%\red{Understand \computationCircuits{} as preparation of q-samples to normalized basis encodings.}

% Orientation on Low construction: Motivation by sparsity
In the previous sections we have studied the schemes to encode arbitrary conditional probability distributions, which concatenation prepares q-samples of Bayesian networks \cite{low_quantum_2014}.
While these schemes do not assume any further structure in the conditional probability distributions, we now investigate how to exploit the structure to achieve sparser circuits.
To be more precise, we focus on deterministic conditional probability distributions, which are basis encodings to functions.
%
To this end, let us first introduce a quantum pendant to encode functions, which has the decomposition by contraction property as we show later.
We will also relate them to basis encodings.

\begin{definition}[\ComputationCircuit{}]
    Given a boolean function $\exformula:\atomstates\rightarrow[2]$ the \computationCircuit{} is the unitary tensor
    \begin{align*}
        \qcbencodingofat{\exformula}{\cheadvariableof{\exformula}{\insymbol},\cheadvariableof{\exformula}{\outsymbol},\shortcatvariables}
        &\quad=
        \sum_{\shortcatindicesin\wcols\exformulaat{\shortcatindices}=1}
        \paulixat{\cheadvariableof{\exformula}{\insymbol},\cheadvariableof{\exformula}{\outsymbol}} \otimes
        \onehotmapofat{\shortcatindices}{\shortcatvariables} \\
        & \quad \quad \quad +
        \sum_{\shortcatindicesin\wcols\exformulaat{\shortcatindices}=0}
        \identityat{\cheadvariableof{\exformula}{\insymbol},\cheadvariableof{\exformula}{\outsymbol}} \otimes
        \onehotmapofat{\shortcatindices}{\shortcatvariables} \, .
    \end{align*}
\end{definition}

\begin{example}[CNOT gate]\label{exa:cnotBenc}
    The CNOT gate is the controlled Pauli-X, that is
    \begin{align*}
        \mathrm{CNOT}\left[\cheadvariable{\insymbol},\cheadvariable{\outsymbol},\catvariable\right]
        = \paulixat{\cheadvariable{\insymbol},\cheadvariable{\outsymbol}} \otimes \tbasisat{\catvariable}
        + \identityat{\cheadvariable{\insymbol},\cheadvariable{\outsymbol}} \otimes \fbasisat{\catvariable} \, .
    \end{align*}
    For the identity function
    \begin{align*}
        \mathrm{Id}(x) = x \, .
    \end{align*}
    we have
    \begin{align*}
        \qcbencodingofat{\mathrm{Id}}{\cheadvariable{\insymbol},\cheadvariable{\outsymbol},\catvariable}
        = \mathrm{CNOT}\left[\cheadvariable{\insymbol},\cheadvariable{\outsymbol},\catvariable\right] \, .
    \end{align*}
    The $\mathrm{CNOT}$ is therefore the \computationCircuit{} to the identity function.
\end{example}

% Multiple output variables
Functions with multiple output variables, i.e. $\exformula:\atomstates\rightarrow\bigtimes_{\selindexin}[2]$, can be encoded image coordinate wise as a concatenation of the respective circuits.

\subsection{Relation with basis encodings}

\BasisEncodings{} are schemes to encode functions into directed boolean tensors, which enable reasoning about functions by tensor network contractions (see \cite{goessmann_tensor_2026}).
To connect with this formalism, let us now show the relation to the \computationCircuits{}.

\begin{lemma}
    \label{lem:bencComCircuit}
    For any boolean function we have (see \figref{fig:bencComCircuit})
    \begin{align*}
        \qcbencodingofat{\exformula}{\cheadvariableof{\exformula}{\insymbol},\cheadvariableof{\exformula}{\outsymbol},\shortcatvariables}
        = \contractionof{
            \bencodingofat{\exformula}{\headvariableof{\exformula},\shortcatvariables},
            \bencodingofat{\oplus}{\cheadvariableof{\exformula}{\outsymbol},\cheadvariableof{\exformula}{\insymbol},\headvariableof{\exformula}}
        }{\cheadvariableof{\exformula}{\insymbol},\cheadvariableof{\exformula}{\outsymbol},\shortcatvariables} \, .
    \end{align*}
\end{lemma}

% Interpretation as q-sample
The \computationCircuit{} applied on the initial state
\begin{align*}
    \sqrt{\frac{1}{2^{\catorder}}} \onesat{\ccatvariableof{[\atomorder]}{\outsymbol}} \otimes \fbasisat{\cheadvariable{\insymbol}}
\end{align*}
prepares a q-sample to the normalized basis encoding. % Called that basis encoding state!


% Sparse decomposition
In \cite{goessmann_tensor_2026} (and more detailed in \cite{goessmann_tensor-network_2025}) sparse decompositions of \basisEncodings{} have been studied.
Based on the equivalence of \computationCircuits{} and contracted \basisEncodings{}, stated by \lemref{lem:bencComCircuit} we transfer these decomposition schemes in the following.

\begin{figure}[t]
    \begin{center}
        \begin{tikzpicture}[scale=0.35, thick] % , baseline = -3.5pt

    %% Statistic Qubits
    \node[anchor=center] (text) at (-1,9) {\colorlabelsize $\cheadvariableof{\exformula}{\insymbol}$};
    \node[anchor=center] (text) at (12,9) {\colorlabelsize $\cheadvariableof{\exformula}{\outsymbol}$};

    \draw[->-] (1,9) -- (4,9);

    \draw (4,8) rectangle (7,10);
    \node[anchor=center] (text) at (5.5,9) {\corelabelsize $\bencodingof{\oplus}$};
    \draw[->-] (7,9) -- (10,9);

    \draw[->-] (5.5,7) -- (5.5,8) node[midway,right] {\colorlabelsize $\headvariableof{\exformula}$};
    \draw (4,5) rectangle (7,7);
    \node[anchor=center] (text) at (5.5,6) {\corelabelsize $\bencodingof{\exformula}$};

    \draw[->-] (5,4) -- (5,5);
    \drawvariabledot{5}{4}
    \draw[->-] (6,1) -- (6,5);
    \drawvariabledot{6}{1}

%    \draw[dashed] (0.5,7.5) -- (13.5,7.5);

    \draw (1,4) -- (10,4);

    \node[anchor=center] (text) at (3,2.75) {$\vdots$};
    \node[anchor=center] (text) at (-1,2.5) {\colorlabelsize $\shortcatvariables$};

    \node[anchor=center] (text) at (7,2.75) {$\vdots$};
    \node[anchor=center] (text) at (12,2.5) {\colorlabelsize $\shortcatvariables$};

    \draw (1,1) -- (10,1);

    \node[anchor=center] (text) at (-5,5) {${=}$};

    \begin{scope}
        [shift={(-20,0)}]

        \node[anchor=center] (text) at (-1,9) {\colorlabelsize $\cheadvariableof{\exformula}{\insymbol}$};
        \node[anchor=center] (text) at (12,9) {\colorlabelsize $\cheadvariableof{\exformula}{\outsymbol}$};

        \draw (5,4) -- (5,8);
        \drawvariabledot{5}{4}
        \draw (6,1) -- (6,8);
        \drawvariabledot{6}{1}

        \draw (1,9) -- (4,9);
        \draw (4,8) rectangle (7,10);
        \node[anchor=center] (text) at (5.5,9) {\corelabelsize $\qcbencodingof{\exformula}$};
        \draw (7,9) -- (10,9);

        \draw (1,4) -- (10,4);

        %\draw[dashed] (0.5,7.5) -- (10.5,7.5);

        \node[anchor=center] (text) at (3,2.75) {$\vdots$};
        \node[anchor=center] (text) at (-1,2.5) {\colorlabelsize $\shortcatvariables$};

        \node[anchor=center] (text) at (8,2.75) {$\vdots$};
        \node[anchor=center] (text) at (12,2.5) {\colorlabelsize $\shortcatvariables$};

        \draw (1,1) -- (10,1);

    \end{scope}

\end{tikzpicture}
    \end{center}
    \caption{Relation of computation circuit and basis encodings.
    The \computationCircuit{} of a function is the contraction of its \basisEncoding{} with the \basisEncoding{} of the mod2 sum $\oplus$.}
    \label{fig:bencComCircuit}
\end{figure}

% Good place?
\begin{example}[Multiple-controlled NOT]\label{exa:mcnotBenc}
    Let us consider the $\catorder$-ary controlled $\mathrm{NOT}$ gate
    \begin{align*}
        \mathrm{MCNOT}^{\catorder}\left[\cheadvariable{\insymbol},\cheadvariable{\outsymbol},\shortcatvariables\right]
        = \paulixat{\cheadvariable{\insymbol},\cheadvariable{\outsymbol}} \otimes \onehotmapofat{\onetuple{[\catorder]}}{\shortcatvariables}
        + \identityat{\cheadvariable{\insymbol},\cheadvariable{\outsymbol}} \otimes \left(\onesat{\shortcatvariables}-\onehotmapofat{\onetuple{[\catorder]}}{\shortcatvariables}\right) \, .
    \end{align*}
    This is the \computationCircuit{} of the $\catorder$-ary logical conjunction
    \begin{align*}
        \land^{\catorder} \defcols [2]^{\catorder} \rightarrow [2]
        \quad , \quad
         \land^{\catorder}(\shortcatindices)
        = \begin{cases}
              1 & \ifspace \shortcatindices = \onetuple{[\catorder]} \\
              0 & \text{else}
        \end{cases} \, .
    \end{align*}
    Special instances are the $\mathrm{CNOT}$ for $\catorder=1$ (see \exaref{exa:cnotBenc}) and the $\mathrm{Toffoli}$ gate for $\catorder=2$.
\end{example}

\subsection{Circuit \decompositionSparsity{}}

We develop now a quantum equivalent to the circuit decompositions of basis encodings (see \secref{sec:circuits}) sis now a composition of circuits property, as stated in the next lemma.
Similar to the introduction of hidden variables, we here utilize further working qubits to realize decompositions.
To highlight the similarity, we call these sparsity principle circuit \decompositionSparsity{}.

\begin{lemma}
    \label{lem:comCirDecSpar}
    We have for functions $\exfunction:\atomstates\rightarrow\bigtimes_{\selindexin}[2]$, $\secexfunction:\bigtimes_{\selindexin}[2]\rightarrow\bigtimes_{s\in[r]}[2]$ (see \figref{fig:qcbencodingDecomposition})
    \begin{align*}
        &\qcbencodingofat{\secexfunction\circ\exfunction}{
            \cheadvariableof{\secexfunction\circ\exfunction}{\insymbol},
            \cheadvariableof{\secexfunction\circ\exfunction}{\outsymbol},
            \shortcatvariables}
        \otimes \fbasisat{\cheadvariableof{\exfunction}{\outsymbol}} \\
        &\quad= \breakablecontractionof{\onehotmapofat{0}{\cheadvariableof{\exfunction}{\insymbol}},
            \qcbencodingofat{\exfunction}{\cheadvariableof{\exfunction}{\insymbol},
                \cheadvariableof{\exfunction}{\midsymbol},
                \shortcatvariables},
            \qcbencodingofat{\secexfunction}{
                \headvariableof{\insymbol,\secexfunction\circ\exfunction},
                \cheadvariableof{\secexfunction\circ\exfunction}{\outsymbol},
                \cheadvariableof{\exfunction}{\midsymbol}},
            \qcbencodingofat{\exfunction}{\cheadvariableof{\exfunction}{\midsymbol},
                \cheadvariableof{\exfunction}{\outsymbol},
                \shortcatvariables}
        }{
            \cheadvariableof{\secexfunction\circ\exfunction}{\insymbol},\cheadvariableof{\secexfunction\circ\exfunction}{\outsymbol},\shortcatvariables,\cheadvariableof{\exfunction}{\outsymbol}
        } \, .
    \end{align*}
    Here each copy of $\headvariableof{\secexfunction\circ\exfunction}$ denotes a tuple of $r$ boolean variables and $\headvariableof{\exfunction}$ of $\seldim$ boolean variables.
\end{lemma}
\begin{proof}
    We show the claim based on basis calculus and sketch the steps in \figref{fig:proofDecomposition}.
    In the first equation we use \lemref{lem:bencComCircuit} and represent each \computationCircuit{} by a contraction of two basis encodings.
    In the second equation we use the identity
    \begin{align*}
        \contractionof{\bencodingofat{\oplus}{\cheadvariable{\outsymbol},\cheadvariable{\insymbol},\headvariableof{\exfunction}},
            \fbasisat{\cheadvariable{\insymbol}}}{\cheadvariable{\outsymbol},\headvariableof{\exfunction}}
        = \identityat{\cheadvariable{\outsymbol},\headvariableof{\exfunction}} \, .
    \end{align*}
    Now in the third equation we use that the contraction of basis encodings is the basis encoding of the composition function, and that $\exfunction\oplus\exfunction$ is the vanishing function, which basis encoding is a tensor product with $\fbasisat{\cheadvariableof{\exfunction}{\outsymbol}}$.
    We use we \lemref{lem:bencComCircuit} again in the fourth equation to arrive at the claim.
%    For any $\shortcatindices\in[2]^{\atomorder}$ we have that (see \figref{fig:proofDecomposition}):
%    \red{Uses the polynomial sparsity to uncompute: $\exfunction \oplus\exfunction = 0$.}
%    Can be seen by the basis encoding representation, when starting the hidden qubit with $\fbasis$, since
%    \begin{align*}
%        \contractionof{\bencodingofat{\oplus}{\headvariable,X_1},\fbasisat{X_0}}{Y,X_1}
%        = \identityat{Y,X_1} \, .
%    \end{align*}
%    The marginalization of the hidden target qubit $X_1$ for measurements does not change the distribution, since we have a deterministic dependence on the measured data register (need preparation by $\fbasis$ for that).
\end{proof}

When having a syntactical decomposition of a propositional formula, we can iteratively apply the \computationCircuit{} decomposition theorem and prepare each connective by a circuit.
We can decompose any propositional formula into logical connectives and prepare to each a \computationCircuit{} (exploiting further sparsity mechanisms).
This works, when the target qubit of one connective is used as a value qubit of another.

%% Closure of
%Note that to exploit \decompositionSparsity{}, we introduce another auxiliary working qubit representing the satisfaction of the sub-function $\exfunction$.
%In \lemref{lem:comCirDecSpar} the equality of the concatenated \computationCircuits{} holds after contraction of this qubit, see also \figref{fig:qcbencodingDecomposition}.
%When performing a computational basis measurement on the state prepared by a concatenation, and omitting the qubit (either not measure it in the first place, or ignore its measurement), one effectively has closed the variable before measurement.
%This holds true due to the deterministic dependence of the measurement outcome on the measurement of the distributed qubits $\shortcatvariables$, but is not true for more generic concatenations of controlled unitaries.
This entailment of the working qubit with the distributed qubits can be resolved by an un-computation circuit, which is the adjoint of the \computationCircuit{} for the working qubits (see \figref{fig:proofDecomposition}).
Note that while this is necessary in some applications, it is not necessary for the quantum rejection sampling presented here.

%\begin{figure}
%    \begin{center}
%        \input{tikz/computation-circuits/decomposition_sparsity}
%    \end{center}
%    \caption{
%        Exploitation of \DecompositionSparsity{} in \computationCircuit{}s.
%        The working qubit $\headvariableof{\exfunction}$ is initialized in the ground state $\fbasis$ and un-computed after its usage to prepare $\headvariableof{\secexfunction\circ\exfunction}$.
%        It is then disentangled with the other qubits.
%    }\label{fig:qcbencodingDecomposition}
%\end{figure}

\begin{figure}
    \begin{center}
        \begin{tikzpicture}[scale=0.35, thick] % , baseline = -3.5pt

    \node[anchor=center] (text) at (2,9) {\colorlabelsize $\cheadvariableof{\secexfunction\circ\exfunction}{\insymbol}$};
    \node[anchor=center] (text) at (14,9) {\colorlabelsize $\cheadvariableof{\secexfunction\circ\exfunction}{\outsymbol}$};

    \draw (3,9) -- (7,9);
    \draw (7,8) rectangle (9,10);
    \node[anchor=center] (text) at (8,9) {\corelabelsize $\qcbencodingof{\secexfunction}$};
    \draw (9,9) -- (13,9);

    \draw (8,6) -- (8,8);
    \drawvariabledot{8}{6}

    \draw (1,5) rectangle (3,7);
    \drawtextnode{2}{6}{\corelabelsize $\onehotmapof{0}$};
    \draw (3,6) -- (5,6);
    \draw (5,5) rectangle (7,7);
    \node[anchor=center] (text) at (6,6) {\corelabelsize $\qcbencodingof{\exfunction}$};
    \draw (7,6) -- (9,6);
    \draw (5.5,4) -- (5.5,5);
    \drawvariabledot{5.5}{4}
    \draw (6.5,1) -- (6.5,5);
    \drawvariabledot{6.5}{1}

    \draw (9,5) rectangle (11,7);
    \node[anchor=center] (text) at (10,6) {\corelabelsize $\qcbencodingof{\exfunction}$};
    \draw (11,6) -- (13,6);
    \drawtextnode{14}{6}{\colorlabelsize $\cheadvariableof{\exfunction}{\outsymbol}$}
    \draw (9.5,4) -- (9.5,5);
    \drawvariabledot{9.5}{4}
    \draw (10.5,1) -- (10.5,5);
    \drawvariabledot{10.5}{1}

    \draw (5.5,4) -- (5.5,5);
    \drawvariabledot{5.5}{4}
    \draw (6.5,1) -- (6.5,5);
    \drawvariabledot{6.5}{1}

    \draw[dashed] (0.5,7.5) -- (13.5,7.5);

    \draw (3,4) -- (13,4);

    %\draw (1,0.5) rectangle (3,4.5);
    %\drawtextnode{2}{2.5}{\corelabelsize $\onehotmapof{\shortcatindices}$}

    \node[anchor=center] (text) at (5,2.75) {$\vdots$};
    \node[anchor=center] (text) at (2,2.5) {\colorlabelsize $\shortcatvariables$};

    \node[anchor=center] (text) at (11,2.75) {$\vdots$};
    \node[anchor=center] (text) at (14,2.5) {\colorlabelsize $\shortcatvariables$};

    \draw (3,1) -- (13,1);

    \drawtextnode{15}{5}{${=}^{(1)}$}

    \begin{scope}[shift={(16,-3)}]

        \begin{scope}[shift={(0,3)}]
            \node[anchor=center] (text) at (2,12) {\colorlabelsize $\cheadvariableof{\secexfunction\circ\exfunction}{\insymbol}$};
            \node[anchor=center] (text) at (14,12) {\colorlabelsize $\cheadvariableof{\secexfunction\circ\exfunction}{\outsymbol}$};
            \draw[->-] (9,12) -- (13,12);
            \draw[->-] (3,12) -- (7,12);

            \draw (7,11) rectangle (9,13);
            \drawtextnode{8}{12}{\corelabelsize $\bencodingof{\oplus}$}
            \draw[->-] (8,10) -- (8,11);
            \draw (7,8) rectangle (9,10);
            \node[anchor=center] (text) at (8,9) {\corelabelsize $\bencodingof{\secexfunction}$};


            \draw[->] (8,6) -- (8,8);
            \drawvariabledot{8}{6}

            \draw (1,5) rectangle (3,7);
            \drawtextnode{2}{6}{\corelabelsize $\onehotmapof{0}$};
            \draw[->-] (3,6) -- (5,6);

            \draw[->-] (11,6) -- (13,6);
            \drawtextnode{14}{6}{\colorlabelsize $\cheadvariableof{\exfunction}{\outsymbol}$}

        \end{scope}
        %% Statistic Qubits

        \draw (5,8) rectangle (7,10);
        \drawtextnode{6}{9}{\corelabelsize $\bencodingof{\oplus}$}
        \draw[->-] (6,7) -- (6,8);
        \draw (5,5) rectangle (7,7);
        \node[anchor=center] (text) at (6,6) {\corelabelsize $\bencodingof{\exfunction}$};
        \draw[->-] (5.5,4) -- (5.5,5);
        \drawvariabledot{5.5}{4}
        \draw[->-] (6.5,1) -- (6.5,5);
        \drawvariabledot{6.5}{1}


        \draw[->-] (7,9) -- (8,9);
        \draw[->-] (8,9) -- (9,9);

        \draw (9,8) rectangle (11,10);
        \drawtextnode{10}{9}{\corelabelsize $\bencodingof{\oplus}$}
        \draw[->-] (10,7) -- (10,8);
        \draw (9,5) rectangle (11,7);
        \node[anchor=center] (text) at (10,6) {\corelabelsize $\bencodingof{\exfunction}$};

        \draw[->-] (9.5,4) -- (9.5,5);
        \drawvariabledot{9.5}{4}
        \draw[->-] (10.5,1) -- (10.5,5);
        \drawvariabledot{10.5}{1}

        \draw (5.5,4) -- (5.5,5);
        \drawvariabledot{5.5}{4}
        \draw (6.5,1) -- (6.5,5);
        \drawvariabledot{6.5}{1}

        %\draw[dashed] (0.5,7.5) -- (13.5,7.5);

        \draw (3,4) -- (13,4);

        \node[anchor=center] (text) at (5,2.75) {$\vdots$};
        \node[anchor=center] (text) at (2,2.5) {\colorlabelsize $\shortcatvariables$};
        %\draw (1,0.5) rectangle (3,4.5);
        %\drawtextnode{2}{2.5}{\corelabelsize $\onehotmapof{\shortcatindices}$}

        \node[anchor=center] (text) at (11,2.75) {$\vdots$};
        \node[anchor=center] (text) at (14,2.5) {\colorlabelsize $\shortcatvariables$};

        \draw (3,1) -- (13,1);
    \end{scope}

    \drawtextnode{33}{5}{${=}^{(2)}$}

    \begin{scope}[shift={(32,-1.5)}]

        \node[anchor=center] (text) at (2,12) {\colorlabelsize $\cheadvariableof{\secexfunction\circ\exfunction}{\insymbol}$};
        \node[anchor=center] (text) at (14,12) {\colorlabelsize $\cheadvariableof{\secexfunction\circ\exfunction}{\outsymbol}$};
        \draw[->-] (7,12) -- (13,12);
        \draw[->-] (3,12) -- (5,12);

        \draw (5,11) rectangle (7,13);
        \drawtextnode{6}{12}{\corelabelsize $\bencodingof{\oplus}$}
        \draw[->-] (6,10) -- (6,11);

        \draw (5,8) rectangle (7,10);
        \drawtextnode{6}{9}{\corelabelsize $\bencodingof{\secexfunction}$}
        \draw[->-] (6,7) -- (6,8);
        \draw (5,5) rectangle (7,7);
        \node[anchor=center] (text) at (6,6) {\corelabelsize $\bencodingof{\exfunction}$};
        \draw[->-] (5.5,4) -- (5.5,5);
        \drawvariabledot{5.5}{4}
        \draw[->-] (6.5,1) -- (6.5,5);
        \drawvariabledot{6.5}{1}

        \draw[->-] (11,10) -- (13,10);
        \drawtextnode{14}{10}{\colorlabelsize $\cheadvariableof{\exfunction}{\outsymbol}$}
        \draw (9,9) rectangle (11,11);
        \drawtextnode{10}{10}{\corelabelsize $\bencodingof{\oplus}$}
        \draw[->-] (10,8) -- (10,9);
        \draw[->-] (10,8) to[bend left=90] (9,10);
        \drawvariabledot{10}{8}
        \draw[->-] (10,7) -- (10,8);
        \draw (9,5) rectangle (11,7);
        \node[anchor=center] (text) at (10,6) {\corelabelsize $\bencodingof{\exfunction}$};

        \draw[->-] (9.5,4) -- (9.5,5);
        \drawvariabledot{9.5}{4}
        \draw[->-] (10.5,1) -- (10.5,5);
        \drawvariabledot{10.5}{1}

        \draw (5.5,4) -- (5.5,5);
        \drawvariabledot{5.5}{4}
        \draw (6.5,1) -- (6.5,5);
        \drawvariabledot{6.5}{1}

        %\draw[dashed] (0.5,7.5) -- (13.5,7.5);

        \draw (3,4) -- (13,4);

        %\draw (1,0.5) rectangle (3,4.5);
        %\drawtextnode{2}{2.5}{\corelabelsize $\onehotmapof{\shortcatindices}$}
        \node[anchor=center] (text) at (5,2.75) {$\vdots$};
        \node[anchor=center] (text) at (2,2.5) {\colorlabelsize $\shortcatvariables$};
        \node[anchor=center] (text) at (11,2.75) {$\vdots$};
        \node[anchor=center] (text) at (14,2.5) {\colorlabelsize $\shortcatvariables$};

        \draw (3,1) -- (13,1);
    \end{scope}

    \begin{scope}[shift={(7,-15)}]

        \drawtextnode{0}{5}{${=}^{(3)}$}

        \node[anchor=center] (text) at (2,9) {\colorlabelsize $\cheadvariableof{\secexfunction\circ\exfunction}{\insymbol}$};
        \node[anchor=center] (text) at (14,9) {\colorlabelsize $\cheadvariableof{\secexfunction\circ\exfunction}{\outsymbol}$};
        \draw[->-] (7,9) -- (13,9);
        \draw[->-] (3,9) -- (5,9);

        \draw (5,8) rectangle (7,10);
        \drawtextnode{6}{9}{\corelabelsize $\bencodingof{\oplus}$}
        \draw[->-] (6,7) -- (6,8);
        \draw (5,5) rectangle (7,7);
        \node[anchor=center] (text) at (6,6) {\corelabelsize $\bencodingof{\secexfunction\circ\exfunction}$};
        \draw[->-] (5.5,4) -- (5.5,5);
        \drawvariabledot{5.5}{4}
        \draw[->-] (6.5,1) -- (6.5,5);
        \drawvariabledot{6.5}{1}

        \draw (11,6) -- (13,6);
        \drawtextnode{14}{6}{\colorlabelsize $\cheadvariableof{\exfunction}{\outsymbol}$}

        \draw (9,5) rectangle (11,7);
        \node[anchor=center] (text) at (10,6) {\corelabelsize $\fbasis$};

        \draw (3,4) -- (13,4);

        %\draw (1,0.5) rectangle (3,4.5);
        %\drawtextnode{2}{2.5}{\corelabelsize $\onehotmapof{\shortcatindices}$}
        \node[anchor=center] (text) at (5,2.75) {$\vdots$};
        \node[anchor=center] (text) at (2,2.5) {\colorlabelsize $\shortcatvariables$};
        \node[anchor=center] (text) at (11,2.75) {$\vdots$};
        \node[anchor=center] (text) at (14,2.5) {\colorlabelsize $\shortcatvariables$};

        \draw (3,1) -- (13,1);
    \end{scope}

    \begin{scope}[shift={(27,-15)}]
        \drawtextnode{-3}{5}{${=}^{(4)}$}

        \node[anchor=center] (text) at (-1,9) {\colorlabelsize $\cheadvariableof{\secexfunction\circ\exfunction}{\insymbol}$};
        \node[anchor=center] (text) at (12,9) {\colorlabelsize $\cheadvariableof{\secexfunction\circ\exfunction}{\outsymbol}$};
        \drawtextnode{12}{6}{\colorlabelsize $\cheadvariableof{\exfunction}{\outsymbol}$}

        \draw (5,4) -- (5,8);
        \drawvariabledot{5}{4}
        \draw (6,1) -- (6,8);
        \drawvariabledot{6}{1}

        \draw (1,9) -- (4,9);
        \draw (4,8) rectangle (7,10);
        \node[anchor=center] (text) at (5.5,9) {\corelabelsize $\qcbencodingof{\secexfunction\circ\exfunction}$};
        \draw (7,9) -- (10,9);

        \draw (1,4) -- (10,4);

        \draw[dashed] (0.5,7.5) -- (10.5,7.5);

        \node[anchor=center] (text) at (3,2.75) {$\vdots$};
        \node[anchor=center] (text) at (-1,2.5) {\colorlabelsize $\shortcatvariables$};

        \node[anchor=center] (text) at (8,2.75) {$\vdots$};
        \node[anchor=center] (text) at (12,2.5) {\colorlabelsize $\shortcatvariables$};

        \draw (1,1) -- (10,1);

        % Working qubit
        \draw (-1,5) rectangle (1,7);
        \drawtextnode{0}{6}{\corelabelsize $\fbasis$};
        \draw (1,6) -- (10,6);
    \end{scope}



\end{tikzpicture}
    \end{center}
    \caption{
        Proof of \lemref{lem:comCirDecSpar} based on basis calculus.
    }\label{fig:proofDecomposition}
\end{figure}

\begin{example}\label{exa:mcnotDecomposition}
    Let us reconsider the $\catorder$-ary controlled $\mathrm{NOT}$ gate from \exaref{exa:mcnotBenc}.
    The $\catorder$-ary logical conjunction is decomposed as a sequence of $2$-ary logical conjunctions.
    The \computationCircuit{} of each $2$-ary conjunction is a $\mathrm{Toffoli}$ gate applied to an input qubit and a worker qubit.
    Using \DecompositionSparsity{} we therefore find the standard decomposition of $\catorder$-ary controlled $\mathrm{NOT}$ gates into $\mathrm{Toffoli}$ gates.
\end{example}

\subsection{Circuit \polynomialSparsity{}}\glsadd{sc:reedMullerExpansion:qc}

While we have introduced additional auxiliary qubits to represent decomposable functions, we now investigate concatenations of \computationCircuits{}, which are sharing the same target qubit.
As we will show, these are related to summations $\modtwoword$ (which are contractions in the semiring $\galoissemiring$, see \secref{sec:semiringGalois}), and can be constructed using Reed-Muller expansions (see \secref{sec:reedMullerExpansion}).
To highlight the similarity to \polynomialSparsity{}, which has been developed in \charef{cha:sparseCalculus}, we call this sparsity principle circuit \polynomialSparsity{}.


%\begin{lemma}
%    \label{lem:comCircuitMod2}
%    The composition of two \computationCircuits{} to functions $\exfunction$ and $\secexfunction$ (see \figref{fig:mod2ComputationCircuit}), where the \computationCircuits{} share the target qubit, is the \computationCircuit{} of their sum mod2.
%    In particular, \computationCircuits{} commute when they have the same target qubit.
%\end{lemma}
%\begin{proof}
%    This follows from the representation of \computationCircuits{} by basis encodings of the corresponding function and the $\modtwoword$ sum $\oplus$ (see \figref{fig:mod2ComputationCircuit}).
%\end{proof}



\begin{lemma}\label{lem:comCircuitMod2}
    Let $\{\formulaofat{\selindex}{\shortcatvariables}\wcols\selindexin\}$ be a set of propositional formulas, and let
    \begin{align*}
        \secexformulaat{\shortcatvariables}
        = \srcontractionofwrt{\sencodingofat{\formula}{\shortcatvariables,\selvariable}}{\shortcatvariables}{\galoissemiring} \,
    \end{align*}
    be their summation $\modtwoword$, which can be expressed using the selection encoding $\sencodingofat{\formula}{\shortcatvariables,\selvariable}$ of the formula set.
    Then we have that
    \begin{align*}
        \contractionof{
            \{\qcbencodingofat{\formulaof{\selindex}}{\cheadvariable{\selindex},\cheadvariable{\selindex+1},\shortcatvariables} \wcols \selindexin\}
        }{\shortcatvariables}
        = \qcbencodingofat{\secexformula}{\cheadvariable{0},\cheadvariable{\seldim},\shortcatvariables} \, .
    \end{align*}
\end{lemma}
\begin{proof}
    We first show the claim for $\seldim=2$, the sum of two formulas $\formulaof{0}$ and $\formulaof{1}$.
    We decompose in \figref{fig:mod2ComputationCircuit} each \computationCircuit{} into a contraction of the functions basis encoding with the basis encoding of $\oplus$, the summation $\modtwoword$.
    By associativity of $\modtwosum$, we can rewrite the contraction using the basis encoding of $\formulaof{0}\modtwosum\formulaof{1}$.

    % Case m=2
    We then show the generic statement by induction over $\seldim$.
    The statement is trivial for $\seldim=1$.
    % Case m>2 -> Could also do a full induction here
    Let us assume that the claim holds for $\seldim-1$ formulas, and let there be $\seldim$ formulas.
    By assumption, the concatenation of the first $\seldim-1$ computation circuits is the computation circuit of
    \begin{align*}
        \bigoplus_{\catenumerator\in[\seldim]}\formulaof{\catenumerator}
    \end{align*}
    Applying the statement for $\seldim=2$, we conclude that the concatenation of all $\seldim$ computation circuits is the computation circuit of
    \begin{align*}
        \secexformula
        =\bigoplus_{\catenumerator\in[\seldim]}\formulaof{\seccatenumerator}
        = \srcontractionofwrt{\sencodingofat{\formula}{\shortcatvariables,\selvariable}}{\shortcatvariables}{\galoissemiring} \, .
    \end{align*}
\end{proof}


\begin{figure}[t]
    \begin{center}
        \begin{tikzpicture}[scale=0.35, xscale = 0.9, thick]

    \draw (5,4) -- (5,8);
    \drawvariabledot{5}{4}
    \draw (6,1) -- (6,8);
    \drawvariabledot{6}{1}

    \draw (9,4) -- (9,8);
    \drawvariabledot{9}{4}
    \draw (10,1) -- (10,8);
    \drawvariabledot{10}{1}

    \draw (2,9) -- (4,9) node[midway,above] {\colorlabelsize $\headvariable^{\insymbol}$};
    \draw (4,8) rectangle (7,10);
    \node[anchor=center] (text) at (5.5,9) {\corelabelsize $\qcbencodingof{\exformula}$};
    \draw (7,9) -- (8,9);
    \draw (8,8) rectangle (11,10);
    \node[anchor=center] (text) at (9.5,9) {\corelabelsize $\qcbencodingof{\secexformula}$};
    \draw (11,9) -- (13,9) node[midway,above] {\colorlabelsize $\headvariable^{\outsymbol}$};

    \draw (2,4) -- (13,4);

    \node[anchor=center] (text) at (3,2.5) {\colorlabelsize $\shortcatvariables$};

    \node[anchor=center] (text) at (8,2.75) {$\vdots$};
    \node[anchor=center] (text) at (12,2.5) {\colorlabelsize $\shortcatvariables$};

    \draw (2,1) -- (13,1);

    \node[anchor=center] (text) at (14.5,5) {${=}$};

    \begin{scope}[shift={(14,0)}]

        \draw[->-] (2,9) -- (4,9) node[midway,above] {\colorlabelsize $\headvariable^{\insymbol}$};

        \draw (4,8) rectangle (7,10);
        \node[anchor=center] (text) at (5.5,9) {\corelabelsize $\bencodingof{\oplus}$};
        \draw[->-] (7,9) -- (8,9);
        \draw (8,8) rectangle (11,10);
        \node[anchor=center] (text) at (9.5,9) {\corelabelsize $\bencodingof{\oplus}$};

        \draw[->-] (5.5,7) -- (5.5,8) node[midway,right] {\colorlabelsize $\headvariable$};
        \draw (4,5) rectangle (7,7);
        \node[anchor=center] (text) at (5.5,6) {\corelabelsize $\bencodingof{\exformula}$};
        \draw[->-] (9.5,7) -- (9.5,8) node[midway,right] {\colorlabelsize $\headvariable$};
        \draw (8,5) rectangle (11,7);
        \node[anchor=center] (text) at (9.5,6) {\corelabelsize $\bencodingof{\secexformula}$};

        \draw[->-] (11,9) -- (13,9) node[midway,above] {\colorlabelsize $\headvariable^{\outsymbol}$};

        \draw[->-] (5,4) -- (5,5);
        \drawvariabledot{5}{4}
        \draw[->-] (6,1) -- (6,5);
        \drawvariabledot{6}{1}

        \draw[->-] (9,4) -- (9,5);
        \drawvariabledot{9}{4}
        \draw[->-] (10,1) -- (10,5);
        \drawvariabledot{10}{1}

        \draw (2,4) -- (13,4);

        \node[anchor=center] (text) at (3,2.5) {\colorlabelsize $\shortcatvariables$};

        \node[anchor=center] (text) at (7,2.75) {$\vdots$};
        \node[anchor=center] (text) at (12,2.5) {\colorlabelsize $\shortcatvariables$};

        \draw (2,1) -- (13,1);

    \end{scope}

    \node[anchor=center] (text) at (28,5) {${=}$};

    \begin{scope}[shift={(27,-1.5)}]

        \draw[->-] (2,12) -- (6,12) node[midway,above] {\colorlabelsize $\headvariable^{\insymbol}$};
        \draw[->-] (9,12) -- (13,12) node[midway,above] {\colorlabelsize $\headvariable^{\outsymbol}$};

        \draw (6,11) rectangle (9,13);
        \node[anchor=center] (text) at (7.5,12) {\corelabelsize $\bencodingof{\oplus}$};
        \draw[->-] (7.5,10) -- (7.5,11);
        \draw (4,8) rectangle (11,10);
        \node[anchor=center] (text) at (7.5,9) {\corelabelsize $\bencodingof{\oplus}$};

        \draw[->-] (5.5,7) -- (5.5,8);% node[midway,right] {\colorlabelsize $\headvariable$};
        \draw (4,5) rectangle (7,7);
        \node[anchor=center] (text) at (5.5,6) {\corelabelsize $\bencodingof{\exformula}$};
        \draw[->-] (9.5,7) -- (9.5,8);% node[midway,right] {\colorlabelsize $\headvariable$};
        \draw (8,5) rectangle (11,7);
        \node[anchor=center] (text) at (9.5,6) {\corelabelsize $\bencodingof{\secexformula}$};



        \draw[->-] (5,4) -- (5,5);
        \drawvariabledot{5}{4}
        \draw[->-] (6,1) -- (6,5);
        \drawvariabledot{6}{1}

        \draw[->-] (9,4) -- (9,5);
        \drawvariabledot{9}{4}
        \draw[->-] (10,1) -- (10,5);
        \drawvariabledot{10}{1}

        \draw (2,4) -- (13,4);

        \node[anchor=center] (text) at (3,2.5) {\colorlabelsize $\shortcatvariables$};

        \node[anchor=center] (text) at (7,2.75) {$\vdots$};
        \node[anchor=center] (text) at (12,2.5) {\colorlabelsize $\shortcatvariables$};

        \draw (2,1) -- (13,1);

    \end{scope}

    \node[anchor=center] (text) at (41,5) {${=}$};

    \begin{scope}[shift={(40,0)}]

        \draw (5,4) -- (5,8);
        \drawvariabledot{5}{4}
        \draw (6,1) -- (6,8);
        \drawvariabledot{6}{1}

        \draw (2,9) -- (4,9) node[midway,above] {\colorlabelsize $\headvariable^{\insymbol}$};
        \draw (4,8) rectangle (7,10);
        \node[anchor=center] (text) at (5.5,9) {\corelabelsize $\qcbencodingof{\exformula\oplus\secexformula}$};
        \draw (7,9) -- (9,9) node[midway,above] {\colorlabelsize $\headvariable^{\outsymbol}$};

        \draw (2,4) -- (9,4);

        \node[anchor=center] (text) at (3,2.5) {\colorlabelsize $\shortcatvariables$};

        \node[anchor=center] (text) at (4,2.75) {$\vdots$};
        \node[anchor=center] (text) at (8,2.5) {\colorlabelsize $\shortcatvariables$};

        \draw (2,1) -- (9,1);

    \end{scope}

\end{tikzpicture}
    \end{center}
    \caption{Equivalence of the composition of \computationCircuits{} with the \computationCircuit{} of their mod2 sum (\lemref{lem:comCircuitMod2}).}
    \label{fig:mod2ComputationCircuit}
\end{figure}


Recall that while the quantum circuit itself is defined in the complex sum-product semiring, we use the semiring $\galoissemiring$ for function decomposition.
% Algebraic Normal Form
The $\cpformatof{\galoissemiring}$ decomposition is closely related to the Algebraic Normal Form of a Boolean polynomial.
In the latter, we need to restrict to leg vectors $\onesat{\catvariable}$ and $\tbasisat{\catvariable}$, that is forbid the usage of $\fbasisat{\catvariable}$ vectors.
The Algebraic Normal Form is also called positive-polarity Reed-Muller expansion.
Allowing for the $\fbasisat{\catvariable}$ vector in the $\cpformatof{\galoissemiring}$ decomposition corresponds with mixed-polarity Reed-Muller (MPRM) expansions.


Let us now determine \computationCircuits{} for MPRM terms to get an encoding for arbitrary formulas in MPRM expansions.

\begin{lemma}\label{lem:mprmTermRep}
    A MPRM term $\formulaof{\hardparam}$ has the \computationCircuit{} (see \figref{fig:qcbencodingPolynomial})
    \begin{align*}
        &\contractionof{\qcaencodingofat{\formulaof{\hardparam}}{\cheadvariable{\insymbol},\cheadvariable{\outsymbol},\shortcatvariables},
            \identityat{\shortcatvariables,\ccatvariableof{[\catorder]}{\insymbol},\ccatvariableof{[\catorder]}{\outsymbol}}}{
            \cheadvariable{\insymbol},\cheadvariable{\outsymbol},\ccatvariableof{[\catorder]}{\insymbol},\ccatvariableof{[\catorder]}{\outsymbol}
        } \\
        &\quad = \breakablecontractionof{
            \{\mathrm{MCNOT}^{\cardof{\hardlegset}}[\cheadvariable{\insymbol},\cheadvariable{\outsymbol},\catvariableof{\hardlegset}]\}
            \cup \left(\bigcup_{\catenumerator\in\hardlegset\ncond\headindexof{\catenumerator}=0}\{\paulixat{\ccatvariableof{\catenumerator}{\insymbol},\catvariableof{\catenumerator}},
                     \paulixat{\catvariableof{\catenumerator},\ccatvariableof{\catenumerator}{\outsymbol}}\}\right) \\
            &\quad\quad\quad\quad\quad \cup \left(\bigcup_{\catenumerator\in\hardlegset\ncond\headindexof{\catenumerator}=1}\{
                                                \identityat{\catvariableof{\catenumerator},\ccatvariableof{\catenumerator}{\insymbol},\ccatvariableof{\catenumerator}{\outsymbol}}
                                                \}\right)
        }{
            \cheadvariable{\insymbol},\cheadvariable{\outsymbol},\ccatvariableof{[\catorder]}{\insymbol},\ccatvariableof{[\catorder]}{\outsymbol}
        } \, .
    \end{align*}
    The \computationCircuit{} of a formula $\formula$ with MPRM expansion $\sliceset$ is a concatenation of the \computationCircuits{} to $\hardlegset\in\sliceset$.
\end{lemma}
\begin{proof}
    For any $\ccatindexof{[\catorder]}{\insymbol},\ccatindexof{[\catorder]}{\outsymbol}$ the slices on both sides are vanishing, if $\ccatindexof{[\catorder]}{\insymbol}\neq\ccatindexof{[\catorder]}{\outsymbol}$.
    If $\ccatindexof{[\catorder]}{\insymbol}=\ccatindexof{[\catorder]}{\outsymbol}$ we have
    \begin{align*}
        &\breakablecontractionof{
            \{\mathrm{MCNOT}^{\cardof{\hardlegset}}[\cheadvariable{\insymbol},\cheadvariable{\outsymbol},\catvariableof{\hardlegset}]\}
            \cup \left(\bigcup_{\catenumerator\in\hardlegset\ncond\headindexof{\catenumerator}=0}\{\paulixat{\ccatvariableof{\catenumerator}{\insymbol},\catvariableof{\catenumerator}},
                     \paulixat{\catvariableof{\catenumerator},\ccatvariableof{\catenumerator}{\outsymbol}}\}\right) \\
            &\quad\quad\quad\quad \cup \left(\bigcup_{\catenumerator\in\hardlegset\ncond\headindexof{\catenumerator}=1}\{
                                           \identityat{\catvariableof{\catenumerator},\ccatvariableof{\catenumerator}{\insymbol},\ccatvariableof{\catenumerator}{\outsymbol}}
                                           \}\right)
        }{
            \cheadvariable{\insymbol},\cheadvariable{\outsymbol},\ccatvariableof{[\catorder]}{\insymbol}=\ccatindexof{[\catorder]}{\insymbol},\ccatvariableof{[\catorder]}{\outsymbol}=\ccatindexof{[\catorder]}{\outsymbol}
        } \\
        & \quad = \contractionof{\{\mathrm{MCNOT}^{\cardof{\hardlegset}}[\cheadvariable{\insymbol},\cheadvariable{\outsymbol},\catvariableof{\hardlegset}]\}
            \cup\{\onehotmapofat{(1-\ccatindexof{\catenumerator}{\insymbol})}{\catvariableof{\catenumerator}} \wcols \catenumerator\in\hardlegset\ncond\headindexof{\catenumerator}=0\}
            \cup\{\onehotmapofat{\ccatindexof{\catenumerator}{\insymbol}}{\catvariableof{\catenumerator}} \wcols \catenumerator\in\hardlegset\ncond\headindexof{\catenumerator}=1\}
        }{ \cheadvariable{\insymbol},\cheadvariable{\outsymbol}}\\
        & \quad =
        \begin{cases}
            \mathrm{NOT}[\cheadvariable{\insymbol},\cheadvariable{\outsymbol}] & \ifspace \formulaof{\hardparam}(\ccatindexof{[\catorder]}{\insymbol})=1 \\
            \identityat{\cheadvariable{\insymbol},\cheadvariable{\outsymbol}} & \ifspace \formulaof{\hardparam}(\ccatindexof{[\catorder]}{\insymbol})=0
        \end{cases} \, .
    \end{align*}
    Thus, the right contraction coincides for all slices on the variables $\ccatvariableof{[\catorder]}{\insymbol}$ and $\ccatvariableof{[\catorder]}{\outsymbol}$ with the oracle of $\formulaof{\hardparam}$.
\end{proof}

\begin{figure}
    \begin{center}
        \begin{tikzpicture}[scale=0.35, thick] % , baseline = -3.5pt

    \node[anchor=center] (text) at (-5,4) {${=}$};

    \begin{scope}
    [shift={(-20,-1)}]

        \node[anchor=center] (text) at (-1,9) {\colorlabelsize $\cheadvariable{\insymbol}$};
        \node[anchor=center] (text) at (12,9) {\colorlabelsize $\cheadvariable{\outsymbol}$};

        \draw (5,6) -- (5,8);
        \drawvariabledot{5}{6}
        \draw (6,-1) -- (6,8);
        \drawvariabledot{6}{-1}

        \draw (1,9) -- (3.5,9);
        \draw (3.5,8) rectangle (7.5,10);
        \node[anchor=center] (text) at (5.5,9) {\corelabelsize $\qcbencodingof{\formulaof{\hardparam}}$};
        \draw (7.5,9) -- (10,9);

        \draw (1,6) -- (10,6);

        \draw[dashed] (0.5,7.5) -- (10.5,7.5);

        \node[anchor=center] (text) at (3,3.25) {$\vdots$};
        \node[anchor=center] (text) at (-1,3) {\colorlabelsize $\ccatvariableof{\variableset}{\insymbol}$};

        \node[anchor=center] (text) at (8,3.25) {$\vdots$};
        \node[anchor=center] (text) at (12,3) {\colorlabelsize $\ccatvariableof{\variableset}{\outsymbol}$};

        \draw (1,-1) -- (10,-1);

    \end{scope}

    \node[anchor=center] (text) at (-1,8) {\colorlabelsize $\cheadvariable{\insymbol}$};
    \node[anchor=center] (text) at (12,8) {\colorlabelsize $\cheadvariable{\outsymbol}$};

    \draw (1,8) -- (10,8);

    %% CNOT
    \drawcnotsymbol{5.5}{8}

    \draw[dashed] (0.5,6.5) -- (10.5,6.5);

    \draw (1,5) -- (2.5,5);
    \draw (2.5,4) rectangle (4.5,6);
    \node[anchor=center] (text) at (3.5,5) {$\paulixsymbol$};
    \draw (4.5,5) -- (6.5,5);
    \draw (6.5,4) rectangle (8.5,6);
    \node[anchor=center] (text) at (7.5,5) {$\paulixsymbol$};
    \draw (8.5,5) -- (10,5);

    \node[anchor=center] (text) at (1.75,3.75) {$\vdots$};
    \node[anchor=center] (text) at (9.25,3.75) {$\vdots$};

    \draw (1,2) -- (2.5,2);
    \draw (2.5,1) rectangle (4.5,3);
    \node[anchor=center] (text) at (3.5,2) {$\paulixsymbol$};
    \draw (4.5,2) -- (6.5,2);
    \draw (6.5,1) rectangle (8.5,3);
    \node[anchor=center] (text) at (7.5,2) {$\paulixsymbol$};
    \draw (8.5,2) -- (10,2);

    \draw (1,0) -- (10,0);
    \draw (1,-2) -- (10,-2);

    \node[anchor=center] (text) at (1.75,-0.75) {$\vdots$};
    \node[anchor=center] (text) at (9.25,-0.75) {$\vdots$};

    \node[anchor=center] (text) at (-1,2) {\colorlabelsize $\ccatvariableof{\variableset}{\insymbol}$};
    \node[anchor=center] (text) at (12,2) {\colorlabelsize $\ccatvariableof{\variableset}{\outsymbol}$};

    \draw (4.75,5) -- (4.75,7.35);
    \draw (5.25,2) -- (5.25,7);
    \draw (5.75,0) -- (5.75,7);
    \draw (6.25,-2) -- (6.25,7.35);

    \drawvariabledot{4.75}{5}
    \drawvariabledot{5.25}{2}
    \drawvariabledot{5.75}{0}
    \drawvariabledot{6.25}{-2}


\end{tikzpicture}
    \end{center}
    \caption{
        \ComputationCircuit{} to a MPRM term $\formulaof{\hardparam}$, consisting in a multiple controlled $\mathrm{NOT}$ gate and Pauli-X gates for legs with $\headindexof{\catenumerator}=0$.
    }\label{fig:qcbencodingPolynomial}
\end{figure}


% Generic encoding
Combining \lemref{lem:mprmTermRep} and \lemref{lem:comCircuitMod2} we can create the \computationCircuit{} of any function based on the $\cpformatof{\galoissemiring}$ decomposition, which is equivalent to Reed-Muller expansions (see \secref{sec:reedMullerExpansion}).


\subsection{Function Encoding}

We investigate how Boolean functions can be encoded in quantum states using \computationCircuits{}.
Depending on the start state, we introduce two function encoding states.

\begin{remark}[Basis and Amplitude Encodings in Quantum Computation \cite{schuld_machine_2021} and \tnreason{}]
    % Compare with Quantum Computing
    Basis Encoding in Quantum Computation refers to the representation of classical $n$ bit strings by $n$ qubit basis states, and is called one-hot encoding in \tnreason{}.
    The Basis Encoding scheme in \tnreason{} goes beyond this scheme and also encodes subsets by sums of one-hot encodings to the members of the set.
    In this way, relations and functions are represented by boolean tensors and contraction of them is refered to as \BasisCalculus{}.

    Amplitude Encoding in Quantum Computation refers to the storage of complex numbers in the amplitudes of quantum states.
    The pendant in \tnreason{} is the Coordinate Encoding scheme, where real numbers are stored in the coordinates of real-valued tensors.
    Compared to Amplitude Encoding, Coordinate Encoding does not have the normalization constraint of quantum states.
    The Amplitude Encoding of the square root of a probability distribution is sometimes called q-sample.
\end{remark}

\subsubsection{Basis encoding}

\begin{definition}
    The basis encoding of a boolean function $\exformula$ is the state
    \begin{align*}
        \basqstateofat{\exformula}{\shortcatvariables,\headvariable}
        = \sqrt{\frac{1}{2^\atomorder}} \sum_{\shortcatindices\in\atomstates} \onehotmapofat{\shortcatindices}{\shortcatvariables}
        \otimes \onehotmapofat{\exformula(\shortcatindices)}{\headvariable} \, .
    \end{align*}
\end{definition}

\begin{lemma}
    When applying the \computationCircuit{} on the initial state
    \begin{align*}
        \sqrt{\frac{1}{2^\atomorder}}\onesat{\shortcatvariables} \otimes \fbasisat{\headvariable} % insymbol?
    \end{align*}
    we get the basis encoding state (see \figref{fig:basSignQstateConstruction}a).
\end{lemma}
\begin{proof}
    Since the Pauli-X turns $\fbasis$ into $\tbasis$, and is applied exactly when $f(\shortcatindices)=1$.
\end{proof}


\begin{corollary}
    Using \lemref{lem:bencComCircuit} we have
    \begin{align*}
        \basqstateofat{\exformula}{\shortcatvariables,\headvariable}
        = \sqrt{\frac{1}{2^{\catorder}}} \bencodingofat{\exformula}{\shortcatvariables,\headvariable}
    \end{align*}
\end{corollary}

\begin{figure}[t]
    \begin{center}
        \begin{tikzpicture}[scale=0.35, thick] % , baseline = -3.5pt

    %% BASIS ENCODING

    \node[anchor=center] (text) at (-6.5,11) {$a)$};

    \begin{scope}[shift={(-5.5,0)}]
        \draw (0,2) rectangle (2,10);
        \node[anchor=center] (text) at (1,6) {\corelabelsize $\basqstateof{\exformula}$};
        \draw (2,9.5) -- (3.5,9.5); % node[anchor=west] {\colorlabelsize $\cheadvariableof{\exformula}{\outsymbol}$};
        \node[anchor=south] (text) at (3.25,9.5) {\colorlabelsize $\cheadvariableof{\exformula}{\outsymbol}$};

        \draw (2,7) -- (3.5,7);
        \node[anchor=center] (text) at (2.75,5.25) {\colorlabelsize $\vdots$};
        \node[anchor=west] (text) at (3,5) {\colorlabelsize $\shortcatvariables$};
        \draw (2,3) -- (3.5,3);
    \end{scope}

    \node[anchor=center] (text) at (0,7) {${=}$};

    \begin{scope}[shift={(4.5,0)}]


        \draw (4,7) -- (4,8.5);
        \drawvariabledot{4}{7}
        \draw (5,3) -- (5,8.5);
        \drawvariabledot{5}{3}

        \draw (-3,8.5) rectangle (-1,10.5);
        \node[anchor=center] (text) at (-2,9.5) {\corelabelsize $\fbasis$};
        \draw (-1,9.5) -- (3,9.5);
        \draw (3,8.5) rectangle (6,10.5);
        \node[anchor=center] (text) at (4.5,9.5) {\corelabelsize $\qcbencodingof{\exformula}$};
        \draw (6,9.5) -- (8,9.5);

        \node[anchor=south] (text) at (7.25,9.5) {\colorlabelsize $\cheadvariableof{\exformula}{\outsymbol}$};

        \draw (-3,8) rectangle (-1,6);
        \node[anchor=center] (text) at (-2,7) {\corelabelsize $\fbasis$};
        \draw (-1,7) -- (0.5,7);
        \draw (0.5,8) rectangle (2.5,6);
        \node[anchor=center] (text) at (1.5,7) {\corelabelsize $\hgate$};
        \draw (2.5,7) -- (8,7);

        \node[anchor=center] (text) at (0,5.25) {\corelabelsize $\vdots$};

        \draw (-3,2) rectangle (-1,4);
        \node[anchor=center] (text) at (-2,3) {\corelabelsize $\fbasis$};
        \draw (-1,3) -- (0.5,3);
        \draw (0.5,2) rectangle (2.5,4);
        \node[anchor=center] (text) at (1.5,3) {\corelabelsize $\hgate$};
        \draw (2.5,3) -- (8,3);

        \node[anchor=center] (text) at (6.5,5.25) {$\vdots$};
        \node[anchor=center] (text) at (8,5) {\colorlabelsize $\shortcatvariables$};

    \end{scope}

    \begin{scope}[shift={(25,0)}]

        %% SIGN ENCODING

        \node[anchor=center] (text) at (-6.5,11) {$b)$};

        \begin{scope}[shift={(-5.5,0)}]
            \draw (0,2) rectangle (2,10);
            \node[anchor=center] (text) at (1,6) {\corelabelsize $\signqstateof{\exformula}$};
            \draw (2,9.5) -- (3.5,9.5); % node[anchor=west] {\colorlabelsize $\cheadvariableof{\exformula}{\outsymbol}$};
            \node[anchor=south] (text) at (3.25,9.5) {\colorlabelsize $\cheadvariableof{\exformula}{\outsymbol}$};

            \draw (2,7) -- (3.5,7);
            \node[anchor=center] (text) at (2.75,5.25) {\colorlabelsize $\vdots$};
            \node[anchor=west] (text) at (3,5) {\colorlabelsize $\shortcatvariables$};
            \draw (2,3) -- (3.5,3);
        \end{scope}

        \node[anchor=center] (text) at (0,7) {${=}$};

        \begin{scope}[shift={(4.5,0)}]


            \draw (4,7) -- (4,8.5);
            \drawvariabledot{4}{7}
            \draw (5,3) -- (5,8.5);
            \drawvariabledot{5}{3}

            \draw (-3,8.5) rectangle (-1,10.5);
            \node[anchor=center] (text) at (-2,9.5) {\corelabelsize $\tbasis$};
            \draw (-1,9.5) -- (0.5,9.5);
            \draw (0.5,8.5) rectangle (2.5,10.5);
            \node[anchor=center] (text) at (1.5,9.5) {\corelabelsize $\hgate$};
            \draw (2.5,9.5) -- (3,9.5);

            \draw (3,8.5) rectangle (6,10.5);
            \node[anchor=center] (text) at (4.5,9.5) {\corelabelsize $\qcbencodingof{\exformula}$};
            \draw (6,9.5) -- (8,9.5);

            \node[anchor=south] (text) at (7.25,9.5) {\colorlabelsize $\cheadvariableof{\exformula}{\outsymbol}$};

            \draw (-3,8) rectangle (-1,6);
            \node[anchor=center] (text) at (-2,7) {\corelabelsize $\fbasis$};
            \draw (-1,7) -- (0.5,7);
            \draw (0.5,8) rectangle (2.5,6);
            \node[anchor=center] (text) at (1.5,7) {\corelabelsize $\hgate$};
            \draw (2.5,7) -- (8,7);

            \node[anchor=center] (text) at (0,5.25) {\corelabelsize $\vdots$};

            \draw (-3,2) rectangle (-1,4);
            \node[anchor=center] (text) at (-2,3) {\corelabelsize $\fbasis$};
            \draw (-1,3) -- (0.5,3);
            \draw (0.5,2) rectangle (2.5,4);
            \node[anchor=center] (text) at (1.5,3) {\corelabelsize $\hgate$};
            \draw (2.5,3) -- (8,3);

            \node[anchor=center] (text) at (6.5,5.25) {$\vdots$};
            \node[anchor=center] (text) at (8,5) {\colorlabelsize $\shortcatvariables$};

        \end{scope}

    \end{scope}

\end{tikzpicture}
    \end{center}
    \caption{Construction of the basis encoding state (a) and the sign encoding state (b) using the \computationCircuit{} of a function.}\label{fig:basSignQstateConstruction}
\end{figure}

The initial state can be prepared by applying Hadamard gates on all distributed variable qubits and preparing the head variable qubit in the ground state.
The overlap of two basis encoding states (which share the same head variable) is proportional to the contraction of the corresponding formulas.

\begin{lemma}
    \label{lem:basQstateOverlap}
    Let $\formula,\secformula$ be two Boolean formulas on the states of $\shortcatvariables$.
    Then we have
    \begin{align*}
        \contraction{\basqstateofat{\formula}{\shortcatvariables,\headvariable},\basqstateofat{\secexformula}{\shortcatvariables,\headvariable}}
        = \frac{1}{2^{\catorder}} \contraction{\formulawith,\secformulawith} \, .
    \end{align*}
\end{lemma}
\begin{proof}
    We have
    \begin{align*}
        \contraction{\basqstateofat{\formula}{\shortcatvariables,\headvariable},\basqstateofat{\secexformula}{\shortcatvariables,\headvariable}}
        &= \frac{1}{2^{\catorder}}
        \sum_{\shortcatindices\in\shortatomstates} \contraction{\onehotmapofat{\formulaat{\indexedshortcatvariables}}{\headvariable},\onehotmapofat{\secformulaat{\indexedshortcatvariables}}{\headvariable}} \\
        &= \frac{1}{2^{\catorder}} \contraction{\formulawith,\secformulawith} \, .
    \end{align*}
\end{proof}

\subsubsection{Sign encoding}

\begin{definition}
    The sign encoding of a boolean function $\exformula$ is the state
    \begin{align*}
        \signqstateofat{\exformula}{\shortcatvariables,\headvariable}
        = \left(\sqrt{\frac{1}{2^\atomorder}} \sum_{\shortcatindices\in\atomstates} \left(-1\right)^{\exformula(\shortcatindices)} \cdot \onehotmapofat{\shortcatindices}{\shortcatvariables}\right)
        \otimes \sqrt{\frac{1}{2}} \left(\fbasisat{\headvariable} - \tbasisat{\headvariable}\right) \, .
    \end{align*}
\end{definition}

Sign encodings can be prepared by the phase kickback trick, as we show next.

\begin{lemma}[Phase kickback trick]\label{lem:phaseKick}
    Applying the \computationCircuit{} $\qcbencodingofat{\formula}{\cheadvariable{\insymbol},\cheadvariable{\outsymbol},\shortcatvariables}$ on a state
    \begin{align*}
        \qstateat{\shortcatvariables} \otimes \hminusbasat{\cavariable{\insymbol}}
    \end{align*}
    we get
    \begin{align*}
        \sum_{\shortcatindicesin} (-1)^{\formula(\shortcatindices)} \qstateat{\indexedshortcatvariables}\cdot \onehotmapofat{\shortcatindices}{\shortcatvariables}
        \otimes \hminusbasat{\cavariable{\outsymbol}} \, .
    \end{align*}
\end{lemma}
\begin{proof}
    We get a disentangled ancilla qubit, because $\hminusbasat{\cavariable{\insymbol}}$ is an eigenvector of the identity and the Pauli-X.
    For the identity, which is applied when $\formula(\shortcatindices)=0$, the eigenvalue is $1$.
    For the Pauli-X, which is applied when $\formula(\shortcatindices)=0$, the eigenvalue is $(-1)$.
\end{proof}

We get the sign encoding with the phase kickback trick, when we initialize with the q-sample of the uniform distribution.

\begin{corollary}% Generalize above: Does not have to be uniform init state !
    When applying the \computationCircuit{} on the initial state
    \begin{align*}
        \sqrt{\frac{1}{2^\atomorder}}\onesat{\shortcatvariables}
        \otimes \sqrt{\frac{1}{2}} \left(\fbasisat{\headvariable} - \tbasisat{\headvariable}\right)
    \end{align*}
    we get the sign encoding state (see \figref{fig:basSignQstateConstruction}b).
    Note that the initial state is a Hadamard gate acting on the state $\tbasisat{\headvariable}$.
\end{corollary}

The sign encoding is also called the phase encoding.
% Other phase encodings
One can ask, whether we can use other ancilla states and get different phase kickbacks while keeping the ancilla disentangled.
However, since we construct here \computationCircuit{} we are restricted to the eigenvectors of the Pauli-X gate, which is the Hadamard basis
\begin{align*}
    \hminusbasat{\headvariable}
    = \sqrt{\frac{1}{2}} \left(\fbasisat{\headvariable} - \tbasisat{\headvariable}\right) \quad,\quad
    \hplusbasat{\headvariable}
    =\sqrt{\frac{1}{2}} \left(\fbasisat{\headvariable} + \tbasisat{\headvariable}\right)
\end{align*}
with eigenvalues $(-1)$ and $0$.


% Preparation of the initial state
The initial state can be prepared by applying Hadamard gates on all distributed variable qubits and preparing the head variable qubit by applying a Hadamard gate on the first basis state, i.e.
\begin{align*}
    \frac{1}{\sqrt{2}} \left( \tbasisat{\headvariableof{\outsymbol}} - \fbasisat{\headvariableof{\outsymbol}} \right)
    = \hgateat{\headvariableof{\outsymbol},\headvariableof{\insymbol}} \fbasisat{\headvariableof{\insymbol}} \, ,
\end{align*}

\begin{lemma}
    \label{lem:signQstateOverlap}
    Let $\formula,\secformula$ be two Boolean formulas on the states of $\shortcatvariables$.
    Then we have
    \begin{align*}
        \contraction{\signqstateofat{\formula}{\shortcatvariables,\headvariable},\signqstateofat{\secexformula}{\shortcatvariables,\headvariable}}
        = \frac{1}{2^{\catorder}} \left(\contraction{\formulawith,\secformulawith}-\contraction{\formulawith,\lnot\secformulawith}\right) \, .
    \end{align*}
\end{lemma}
\begin{proof}
    We have
    \begin{align*}
        \contraction{\signqstateofat{\formula}{\shortcatvariables,\headvariable},\signqstateofat{\secexformula}{\shortcatvariables,\headvariable}}
        &= \frac{1}{2^{\catorder}}
        \sum_{\shortcatindices\in\shortatomstates} (-1)^{\formulaat{\indexedshortcatvariables}+\secformulaat{\indexedshortcatvariables}}  \\
        &= \frac{1}{2^{\catorder}}
        \sum_{\shortcatindices\in\shortatomstates \wcols \formulaat{\indexedshortcatvariables}=\secformulaat{\indexedshortcatvariables}} 1
        + \sum_{\shortcatindices\in\shortatomstates \wcols \formulaat{\indexedshortcatvariables}\neq\secformulaat{\indexedshortcatvariables}} (-1) \\
        &= \frac{1}{2^{\catorder}} \left(\contraction{\formulawith,\secformulawith}-\contraction{\formulawith,\lnot\secformulawith}\right) \, .
    \end{align*}
\end{proof}

Limitations of sign encodings:
\begin{itemize}
    \item The prepared state $\signqstateof{\exformula}$ is distinguished from the state $\signqstateof{\lnot\exformula}$ only by a non-observable phase factor
    \begin{align*}
        \signqstateof{\exformula}
        = (-1) \cdot \signqstateof{\lnot\exformula}
    \end{align*}
    \item With this construction the statistic qubit is disentangled with the distributed qubits (in contrast to basis encoding states).
    Further manipulation of the statistic qubits alone (as we do with \activationCircuits{}) will not retrieve any information about the encoded function.
    \item When doing the sign encoding procedure for multiple formulas, we get a disentangled state of the head variables with the distributed variables encoded in the sign encoding state of $\bigoplus_{\exformulain}\exformula$ (where by $\bigoplus$ we denote the mod2 sum of boolean functions).
\end{itemize}